\section{Literal closed images and a normal target with vertical components}
\textit{Supports: G331, G345, n1, n3, n4.}

For a noninjective finite map factor through $A=R/\ker\to B$ as in n1. Restriction of ideals, uniform limits and actual decreasing representations gives maps $C_R\to C_A$, $D_R\to D_A$. It does not give surjectivity. If $E_R\phi$ is positive, its restriction is an intrinsic competitor, so
\[
(E_R\phi)|_{X_A}\le E_A(\phi|_{X_A}).                 \tag{17}
\]
The positivity hypothesis is needed for this reasoning; arbitrary nonpositive regularizations are not competitors. No ambient extension or restriction equality follows from a closed boundary embedding.

An all-field counterexample has positive CONTINUOUS envelopes at both ends. Put $R=k[x,y]_{(x,y)}$, $Q=R/(xy)$, and gauge $\mathfrak m$. The two curve points are $w_x(x)=1,w_x(y)=\infty$ and $w_y(x)=\infty,w_y(y)=1$. Let $f=-1/v(x)$, continuous in $[-1,0]$. On the regular surface the ordinary order $o$ satisfies $u(v)\le u(o)$ for every $C,D$ function, by the same primary-ideal argument as n4. Since $f(o)=-1$, every competitor is at most $-1$ everywhere; constant $-1$ is a competitor. Thus $p_Rf=E_Rf=-1\in C_R$.

Downstairs on $Q$, the genuine proper primary profiles $-v((x^N,y))/N$ equal $(-1,-1/N)$ and converge uniformly to $(-1,0)$. Thus the restriction of $f$ is itself in $C_Q$, and
\[
E_Q(f|_{X_Q})=(-1,0).
\]
The gap in (17) at $w_y$ is exactly ONE. No $D_R$ extension of this tuple exists: G331's normal zero dichotomy says a $D$ function with value zero at any point on a normal local branch is identically zero, contradicting its value $-1$ at $w_x$. This proves an actual extension obstruction, not just different ideal presentations. The ambient lifts of those quotient profiles increase to the function $0$ where $v(y)<\infty$ and $-1$ where $v(y)=\infty$. It is not USC, since monomial $(v(x),v(y))=(1,M)$ tends to $w_x$. No increasing closure is available.

The intrinsic arities must stay distinct. G345 gives $\mu_R(-1)=\delta_o$ (surface, one input) and $\mu_Q()=\delta_{w_x}+\delta_{w_y}$ (curve, zero inputs). As a finite $R$-module $Q$ has zero ambient degree-two contribution, but that does not erase its intrinsic reference measure or boundary points.

Finally take the NORMAL target $R$ above and the finite reduced INJECTION
\[
R\longrightarrow B=R\times Q,\qquad r\longmapsto(r,\bar r).
\]
Its source envelope for the pulled obstacle is the actual $C_B$ tuple $(-1,(-1,0))$, whereas $E_Rf=-1$. Both are positive continuous, with gap ONE at the vertical lift $w_y$. The source point $w_y$ is isolated on the clopen finite curve section, but its image is not open in $X_R$: monomial values $(M,1)$ approach it. Hence the boundary map is not open. The full fibre maximum is zero at that base point and $-1$ elsewhere; it is not continuous and cannot be in $D_R$ by normal zero dichotomy. The dominating normal-domain trace ignores the vertical generic weight but consequently lacks full fibre support. Normality ALONE cannot replace the no-vertical hypothesis of n1--n3. This example and n4 retain all support and lower-dimensional points and prevent an invented universal compatibility theorem.
